\documentclass[a4paper,10pt]{report}

\def\packagepath{../../preambule}   % path du package princial
\usepackage{\packagepath/preambule}    % utilisation du fichier de configuration

\def\level{Première }              % Classe
\def\course{NSI}              % Matière
\def\eval{DS05 -- Correction}

\def\visibleornot{visible}    % visible or invisible
\def\documentpath{./Sources_Latex}
\def\date{Jeudi 22 Novembre 2025}
\renewcommand{\arraystretch}{1}  % Ecart dans les tableau

\def\ptsexoA{4}
\def\ptsexoB{2}
\def\ptsexoC{2}
\def\ptsexoD{2}
\def\ptsexoE{1}
\def\ptsexoF{4}
\def\ptsexoG{2}
\def\ptsexoH{1}
\def\ptsexoI{2}
\def\ptsexoJ{2}
\def\ptsexoK{1}

\def\ptstotal{\ptsexoA+\ptsexoB}

\usepackage{hyperref}
\usepackage{tasks}

\usetikzlibrary{shapes, arrows, positioning}
\usetikzlibrary{shapes.geometric, arrows, positioning, decorations.pathreplacing}

\tikzstyle{etat} = [draw, rounded corners=15pt, minimum width=2.5cm, minimum height=1.2cm, text centered, font=\bfseries]
\tikzstyle{nouveau} = [etat, fill=blue!40]
\tikzstyle{pret} = [etat, fill=green!60]
\tikzstyle{elu} = [etat, fill=yellow!60]
\tikzstyle{bloque} = [etat, fill=red!60]
\tikzstyle{termine} = [etat, fill=white]
\tikzstyle{fleche} = [thick, ->, >=stealth]

\usepackage{float} % Permet d'utiliser [H]
\usepackage[T1]{fontenc}
\begin{document}

%%%%%%%%%%%%%%%%%%%%%%%%%%%%%%%%%%%%%%%%%%%%%%%%%%ù

%\PageGardeSujetBac[Matiere = Numérique et Sciences Informatiques,
%                  Session = 2025,
%                  AffJour=false,
%                  Duree = {3 heures 30},
%                  DernierePage = \pageref{LastPage},
%                  ModeExamen=false
%                  ]
\renewcommand{\labelitemi}{\textbullet} %pour éviter les tirets dans les "itemize" qui apportent confusion avec le signe moins.
%% Début page de garde style BAC

   %%%%%%%%%%%%%%%%%%%%%%%%%%%%%%%%%%%%%%%%%%%%%%%%%%%%%%%%
%\PageGardeSujetBac[clés]

\pagestyle{DS_FP}          % Feuille de style fancy
\NomPrenomNote{}

\emph{Dans tout le devoir, sauf indication contraire, on considèrera que les nombres sont des entiers relatifs codés sur 8 bits exactement. On supposera que les variables d'entrée de chaque algorithme sont conformes aux attentes de l'algorithme.
}

\exods{\ptsexoA}

\begin{enumerate}
    \item Sans aucune justification, donner le nombre d'états possibles codables avec:
    
\begin{enumerate}[a)]
    \item 3 bits \qquad $2^3=8$
    \item 8 bits \qquad $2^8=256$
    \item n bits \qquad $2^n$
\end{enumerate}

\item On code maintenant un nombre entier relatif avec la méthode du complément à 2 vue en classe. Quelle est la valeur minimale et maximale décimale codable avec:

\begin{enumerate}[a)]
    \item 3 bits \qquad entre $-2^2=-4$ et $2^2-1=3$
    \item 8 bits \qquad entre $-2^7=-128$ et $2^7-1=127$
    \item n bits \qquad entre $-2^{n-1}$ et $2^{n-1}-1$
\end{enumerate}

\item Que peut-on dire de l'entier relatif codé sur 8 bits en complément à deux suivant: \verb|1001 0101|. Sa valeur décimale est-elle: (plusieurs réponses possibles)

\begin{tasks}(4)
    \task[$\bigcirc$] paire
    \task[$\bigotimes$] impaire
    \task[$\bigcirc$] positive
    \task[$\bigotimes$] négative
\end{tasks}

\end{enumerate}



\begin{minipage}{0.49\linewidth}
    \exods{\ptsexoB}

 \'Ecrire l'algorithme de conversion d'un entier naturel compris entre 0 et 255 en binaire sur 8 bits en pseudo code ou en Python.

 \begin{answer}{1}{3.5} 
    \begin{verbatim}
def conversion_dec_to_bin(val_dec):
    binaire = ''
    for i in range(8):
        reste = val_dec % 2
        binaire = str(reste) + binaire
        val_dec = val_dec // 2
    return binaire
    \end{verbatim}   
\end{answer}
\end{minipage}\hfill{}
\begin{minipage}{0.49\linewidth}
    \exods{\ptsexoC}

    Citez les différentes étapes de la conversion d'un entier relatif en binaire sur 8 bits. Vous choisirez de développer un des deux algorithmes vus en cours. Inutile de l'écrire en Python.
    
    \begin{answer}{1}{3.5} 
    \begin{verbatim}
    \end{verbatim}   
\end{answer}
\end{minipage}

\begin{answer}{1}{10}
    \begin{itemize}
    \item Si l'entier est positif, on le convertit en binaire sur 8 bits avec la méthode de la question 1.
    \item Si l'entier est négatif,
    
     Méthode 1:
            \begin{itemize}
                \item On prend la valeur absolue de l'entier.
                \item On le convertit en binaire sur 8 bits avec la méthode de la question 1.
                \item On inverse tous les bits du résultat.
                \item On additionne 1 au résultat précédent.
                \item On obtient le complément à 2 de l'entier.
            \end{itemize}

     Méthode 2:
            \begin{itemize}
                \item On prend la valeur absolue de l'entier.
                \item On le convertit en binaire sur 8 bits avec la méthode de la question 1.
                \item On recopie les bits de la droite vers la gauche jusqu'au premier '1' inclus
                \item On inverse tous les bits à gauche du premier '1' inclus.
                \item On obtient le complément à 2 de l'entier.
            \end{itemize}
    \end{itemize}
    \end{answer}

\medskip



\begin{minipage}{0.49\linewidth}
 \exods{\ptsexoD}

Convertir le nombre entier relatif suivant avec la méthode du complément à 2 sur 8 bits. Justifier.
\begin{answer}{1}{5} 
    \begin{center}
        \verb|-56|
    \end{center}   

         $|-56|=56=$ \verb|0011 1000_2|
        
        Méthode 1:
        \begin{itemize}
            \item Inversion des bits: \verb|1100 0111_2|
            \item Ajouter 1: \verb|1100 1000_2|
        \end{itemize}

        Méthode 2:
        \begin{itemize}
            \item Recopie des bits jusqu'au 1er '1' iclus
            \item Inversion des bits restants \verb|1100 1000_2|
        \end{itemize}
\end{answer}

\end{minipage}\hfill
\begin{minipage}{0.49\linewidth}

     \exods{\ptsexoE}

Convertir le nombre entier relatif suivant avec la méthode du complément à 2 sur 8 bits. Justifer.
\begin{answer}{1}{5} 
    \begin{center}
        \verb|112|
    \end{center}   

    $112$ est un nombre positif donc il se conevertit comme les entiers naturels.

    $112=64+32+16=0111\,0000_2$
\end{answer}
\end{minipage}

\medskip

\pagebreak
\thispagestyle{DS_LP}
\exods{\ptsexoF}

On suppose que l'on dispose des fonctions suivantes en Python (il est inutile d'écrire le code de ces fonctions). Ecrire la fonction \texttt{conversion\_relatif\_binaire} qui convertit un entier relatif en binaire sur 8 bits. Vous utiliserez les fonctions précédentes pour créer votre fonction. 
 Vous appliquerez la méthode de votre choix (la méthode 1 est plus simple !)




  
    \begin{answer}{1}{14.3}
    \begin{verbatim}

def conversion_relatif_binaire(val_entiere):                # METHODE 1
    if val_entiere >= 0:
        return conversion_entier_binaire(val_entiere)
    else:
        val_entiere =  - val_entiere
        binaire = conversion_entier_binaire(val_entiere)

        binaire_inverse = ''
        for bit in binaire:
            binaire_inverse = binaire_inverse + inverse_bit(bit)
        return additionne_binaire(binaire_inverse)

def conversion_relatif_binaire(val_entiere):                # METHODE 2
    if val_entiere >= 0:
        return conversion_entier_binaire(val_entiere)
    else:
        val_entiere =  - val_entiere
        binaire = conversion_entier_binaire(val_entiere)

        resultat = ''
        i = len(binaire) - 1
        premier_1_trouve = False
        # On parcourt le binaire de droite à gauche
        while i >= 0:
            bit = binaire[i]
            
            if not premier_1_trouve:
                resultat = bit + resultat
                if bit == '1':
                    premier_1_trouve = True
            else:
                resultat = inverse_bit(bit) + resultat
            i = i - 1
    
    \end{verbatim}
    \end{answer}


\end{document}